% Part: methods
% Chapter: induction
% Section: introduction

\documentclass[../../../include/open-logic-section]{subfiles}

\begin{document}

\olfileid[tr]{mth}{ind}{int}

\olsection{Giriş}

Tümevarım; mantığın, kuramsal bilgisayar biliminin ve matematiğin hemen
hemen bütün alanlarında farklı biçimlerde kullanılan önemli bir kanıt
tekniğidir. Mantıktaki birçok sonucu kanıtlamak için gereklidir.

Tümevarım çoğu kez tümdengelimle karşılaştırılır ve özelden genele
çıkarım olarak nitelenir. Örneğin çok sayıda yeşil zümrüt gözlemler,
zümrüt diyeceğimiz ama yeşil olmayan hiçbir şey gözlemlemezsek, bütün
zümrütlerin yeşil olduğu sonucuna varabiliriz. Bu, birçok özel durumdan
(bu zümrüt yeşildir, şu zümrüt yeşildir vb.) genel bir iddiaya (bütün
zümrütler yeşildir) ilerlemesi bakımından tümevarımsal bir çıkarımdır.
\emph{Matematiksel} tümevarım da genel bir iddiayla sonuçlanan bir
çıkarımdır; fakat bu ``basit tümevarım''dan çok farklı türdedir.

Kabaca söylersek matematikte tümevarımsal bir kanıt, belirli türdeki
bütün matematiksel nesnelerin belirli bir özelliğe sahip olduğu sonucuna
varır. En basit durumda tümevarımsal kanıtın ele aldığı matematiksel
nesneler doğal sayılardır. Böyle bir durumda tümevarımsal kanıt, bütün
doğal sayıların bir özelliğe sahip olduğunu belirlemek için kullanılır
ve bunu şunları göstererek yapar:
\begin{enumerate}
    \item $0$ bu özelliğe sahiptir ve
    \item bir~$k$ sayısı bu özelliğe sahip olduğunda~$k+1$ de sahiptir.
\end{enumerate}
Doğal sayılar üzerindeki tümevarım, sayı atanabilen matematiksel
nesneler hakkındaki genel iddiaları kanıtlamak için de çoğu kez
kullanılabilir. Örneğin her sonlu kümenin sonlu sayıda~$n$ elemanı
vardır; tümevarımla her~$n$ sayısının ``~$n$ büyüklüğündeki bütün sonlu
kümeler \dots'tır'' özelliğine sahip olduğunu gösterebilirsek, bütün
sonlu kümeler hakkında bir şey göstermiş oluruz.

Tümevarım, \emph{tümevarımsal olarak tanımlanmış} matematiksel nesnelere
de genellenebilir. Örneğin birinci dereceden mantığın ifadeleri gibi
biçimsel bir dilin ifadeleri tümevarımsal olarak tanımlanır.
\emph{Yapısal tümevarım}, bu tür bütün ifadeler hakkındaki sonuçları
kanıtlamanın bir yoludur. Özellikle yapısal tümevarım, mantıkta çok
yararlıdır ve yaygın olarak kullanılır.

\end{document}
